% ===================== RESUMEN EJECUTIVO =====================
\newpage
\section*{Resumen Ejecutivo}
\addcontentsline{toc}{section}{Resumen Ejecutivo}

Se ha realizado el balance de materia y energía de una caldera acuotubular bagacera diseñada para producir \textbf{100 t/h de vapor sobrecalentado} a \SI{106}{barg} y \SI{545}{\celsius}, con agua de alimentación a \SI{270}{\celsius} y una purga continua del 2\%. El combustible es bagazo de caña de azúcar con 48\% de humedad y 10\% de cenizas (AR), representando un escenario conservador de diseño que contempla la materia extraña mineral incorporada durante la cosecha mecanizada. La eficiencia térmica de la caldera es del 94\%.

\textbf{Resultado principal:} Para producir las 100~t/h de vapor se requieren \textbf{37.67 t/h de bagazo}, lo que arroja un ratio de \textbf{2.655 toneladas de vapor por tonelada de bagazo}.

Las propiedades termodinámicas del agua/vapor se calcularon con las tablas IAPWS-97, y el poder calorífico del bagazo se determinó a partir de la fórmula de Hugot. Se incluye además una guía detallada para reproducir estos cálculos mediante simulación en Aspen Plus.

\newpage
